\subsection{18. Deep Learning of Atomically Resolved Scanning Transmission Electron Microscopy Images: Chemical Identification and Tracking Local Transformations}
\textbf{OSTI:} \texttt{1427646}\quad\textbf{Authors:} Ziatdinov, Dyck, Maksov et al. (2017)\quad\textbf{Domain:} Materials / ML / electron microscopy\\
\scorebadge{Coverage}{8}\quad\scorebadge{Agreement}{8}\quad\textbf{Total:} 16/20\par\smallskip
\noindent\textbf{Coverage rationale.} \textbf{Wave~2:} multislice-quality synthetic training set (perovskite [001] with $Z^{1.7}$ HAADF intensity, Poisson noise, scan-line jitter, interface diffuseness, random defects); 7.8M-param U-Net trained in 95~s on A100; 98.8\% pixel accuracy; sub-pixel column RMSE; confusion matrix and peak-detection figures produced. Residual gap: no real experimental STEM frames, no full Prismatic/abTEM multislice.\par
\noindent\textbf{Agreement rationale.} Per-class F1 0.949--0.998, column-centre RMSE sub-pixel (0.06--0.47~px), detection recall $\geq 0.98$ for all atom classes. Consistent with paper's claim that simulation-trained FCN produces quantitative chemical mapping. Direct comparison to paper's real-STEM metrics impossible (no experimental data).\par\smallskip
\noindent\textbf{What reproduced:}
\begin{itemize}[leftmargin=1.4em,itemsep=0.2em,topsep=0.2em]
\item U-Net architecture (5 encoder stages, skip connections, softmax head)
\item AdamW + cosine LR + cross-entropy with class weighting
\item On-the-fly flip/rotate/brightness augmentation
\item Training-from-simulation paradigm
\item Per-pixel 5-class perovskite chemical identity
\item Sub-pixel column-centre detection (RMSE 0.06-0.47 px)
\item Confusion matrix and peak-finding analysis
\end{itemize}\noindent\textbf{What missing:}
\begin{itemize}[leftmargin=1.4em,itemsep=0.2em,topsep=0.2em]
\item Full multislice simulation (Prismatic / abTEM) as training data
\item Evaluation on real experimental HAADF-STEM frames
\item Transformation-tracking (antisite / vacancy dynamics)
\item Time-series analysis of local transformations
\item Noise / dose-reduction sensitivity sweep
\item Generalisation across perovskite chemistries not seen in training
\end{itemize}\noindent\textbf{Follow-on research questions:}
\begin{enumerate}[leftmargin=1.6em,itemsep=0.2em,topsep=0.2em,label=Q\arabic*.]
\item When the synthetic training set is replaced with a small-angle-multislice abTEM simulation (3-4 slices, \textasciitilde{}10x cheaper than full Prismatic), does pixel accuracy hold and does generalisation to real STEM frames improve measurably
\item How does the U-Net's per-class F1 degrade on real experimental HAADF-STEM frames of SrTiO3/LSMO as dose is reduced from 1e5 to 1e3 e-/A\textasciicircum{}2, and at what dose does Sr vs La/Sr confusion become the dominant error
\item Can a self-supervised pretraining stage (masked autoencoder on unlabelled STEM frames) close the synthetic-to-real domain gap with \textless{}5\% real-frame labels
\item Does replacing the U-Net with a Vision Transformer + UperNet head preserve sub-pixel detection precision while providing calibrated classification uncertainty per atom column
\item When the model is applied across a temporal sequence of experimental STEM frames, does a simple Kalman smoother on per-column class probabilities recover transformation trajectories (antisite \textless{}-\textgreater{} vacancy) with known drift statistics
\end{enumerate}

